\clearpage
\renewcommand{\BookRunningTitle}{Inhalt.}
\SourceLeaf{537}\pdfbookmark[0]{Inhalt.}{historical-contents}
\begin{center}\Large INHALT.\par\medskip
\large GAUSS WERKE BAND II. HÖHERE ARITHMETIK.
\end{center}
\BookRule{.25\textwidth}
\par\medskip\Needspace{5\baselineskip}\textit{Abhandlungen.}
\begingroup\fontsize{10}{13}\selectfont
\begin{longtable}{@{}p{.60\linewidth}@{\quad}p{.24\linewidth}@{\quad}r@{}}
Theorematis arithmetici demonstratio nova .\,.\,.\,.\,.\,.\,.\,.\,.\,.\,.\,. & 1808 Jan. .\,. & Seite 1 \\[3pt]
Summatio quarumdam serierum singularium .\,.\,.\,.\,.\,.\,.\,.\,.\,.\,. & 1808 Aug. .\,. & — 9 \\[3pt]
Theorematis fundamentalis in doctrina de residuis quadraticis demonstrationes et ampliationes novae .\,.\,.\,.\,.\,.\,.\,.\,.\,.\,.\,.\,.\,.\,.\,.\,. & 1817 Febr. . & — 47 \\[3pt]
Theoria residuorum biquadraticorum. Commentatio prima .\kern.1em.\kern.4em.\,.\,.\,.\,.\,. & 1825 Apr. . & — 65 \\[3pt]
Theoria residuorum biquadraticorum. Commentatio secunda .\,.\,.\,.\,.\,. & 1831 Apr. . & — 93 \\[3pt]
\end{longtable}\endgroup
\par\medskip\Needspace{5\baselineskip}\textit{Anzeigen eigner Schriften.}
\begingroup\fontsize{10}{13}\selectfont
\begin{longtable}{@{}p{.60\linewidth}@{\quad}p{.24\linewidth}@{\quad}r@{}}
Theorematis arithmetici demonstratio nova .\,.\,.\,.\,.\,.\,.\,.\,.\,.\,.\,. & 1808 Mai .\,. & — 151 \\[3pt]
Summatio quarumdam serierum singularium .\,.\,.\,.\,.\,.\,.\,.\,.\,.\,. & 1808 Sept. . & — 155 \\[3pt]
Theorematis fundamentalis in doctrina de residuis etc. .\,.\,.\,.\,.\,.\,.\,. & 1817 März . & — 159 \\[3pt]
Theoria residuorum biquadraticorum. Comm. I. .\,.\,.\,.\,.\,.\,.\,.\,.\,. & 1825 April . & — 165 \\[3pt]
Theoria residuorum biquadraticorum. Comm. II. .\,.\,.\,.\,.\,.\,.\,.\,.\,. & 1831 April . & — 169 \\[3pt]
\end{longtable}\endgroup
\par\medskip\Needspace{5\baselineskip}\textit{Anzeigen nicht eigner Schriften.}
\begingroup\fontsize{10}{13}\selectfont
\begin{longtable}{@{}p{.60\linewidth}@{\quad}p{.24\linewidth}@{\quad}r@{}}
[\textsc{Dalberg}] Recherches sur l’irréductibilité Arithmétique et Géométrique des nombres et de leurs puissances .\,.\,.\,.\,.\,.\,.\,.\,.\,.\,.\,.\,.\,.\,.\,. & 1809 März . & — 181 \\[3pt]
\textsc{Chernac}. Cribrum Arithmeticum .\,.\,.\,.\,.\,.\,.\,.\,.\,.\,.\,.\,.\,.\,. & 1812 März . & — 181 \\[3pt]
\textsc{Burckhardt}. Tables des diviseurs .\,.\,.\,.\,.\,.\,. & 1814 Nov. 1816 Nov. 1817 Aug. . & — 183 \\[3pt]
\textsc{Erchinger}. Construction des Siebenzehnecks .\,.\,.\,.\,.\,.\,.\,.\,.\,.\,. & 1825 Dec. . & — 186 \\[3pt]
\textsc{Seeber}. Untersuchungen über die Eigenschaften der positiven ternären quadratischen Formen .\,.\,.\,.\,.\,.\,.\,.\,.\,.\,.\,.\,.\,.\,.\,.\,.\,. & 1831 Juli .\,. & — 188 \\[3pt]
\end{longtable}\endgroup
